% LaTeX template for course notes
% Compile with pdfLaTeX

\documentclass[11pt,letter]{ejlnotes}

%\usepackage{cancel}
\usepackage{mathtools}

% Set the header fields
\renewcommand{\coursename}{ENME 570/670: Aerodynamics}
\renewcommand{\lecttopic}{Lifting-Line Theory}
\renewcommand{\authorname}{Eric J. Limacher}
%\renewcommand{\datemodified}{\today}
\renewcommand{\datemodified}{\today\ at \currenttime}

\newcommand{\ts}{\tensym{S}}
\newcommand{\tc}{\tensym{C}}
\newcommand{\va}{\bm{a}}
\newcommand{\vb}{\bm{b}}
\newcommand{\vu}{\bm{u}}
\newcommand{\rl}{\mathbb{R}}
%\newcommand{\n}{\bm{n}}
\newcommand{\he}{\hat{e}}

\usepackage{elmath}

%--------------------
% Document begins
%--------------------
\begin{document}
	% Apply page style to first page
	\thispagestyle{fancy}
	
	% Optional title block
	% Uncomment if you want a centered title on first page:
	% {\centering {\titlefont \lecttopic} \par}
	
	{\bodyfont  % Switch to body font for content
		% Your content goes here

		
		\section*{Overview}

		A standard lifting line algorithm is developed herein using normalized variables.
		
		
		%%%%%%%%%%%%
		\section*{Variable Definitions}
		
		\begin{itemize}
			\item $y = $ span-wise coordinate.
			\item $c = c(y) = $ chord distribution.
			\item $b = $ span.
			\item $\theta = \theta(y) = $ twist distribution.
			\item $\Theta = $ "collective pitch angle" or "wing angle of attack".
			\item $\Gamma = \Gamma(y) = $ circulation distribution. 
			\item $v = v(y) = $ downwash.
			\item $U = $ freestream velocity.
			\item $\alpha = \alpha(y) = $ local effective angle of attack.
			\item $\alpha_0 = $ zero-lift angle of attack for chosen airfoil.
		\end{itemize}

		
		%%%%%%%%%%%%%%%
		\section*{Lifting-Line Theory}
		
		Lifting-line theory represents a straight wing as a line over which circulation is distributed.  The lift per unit length at each location is assumed to obey the Kutta-Joukowsky theorem.  The circulation at each location depends on the local effective angle of attack, which in turn depends on the local twist angle and the downwash induced by the vortex lines trailing behind the wing.  The strength of these trailing vortex lines depends on the spanwise gradient of circulation, $d\Gamma/dy$, as necessary to satisfy Helmholtz's vortex laws.  The circulation distribution, $\Gamma(y)$, must be solved iteratively.  Once $\Gamma(y)$ is known, one can calculate lift and induced drag by integrating the Kutta-Joukowsky theorem along the wing.
		
		To estimate the circulation distribution on wings of arbitrary design, we will discretize the wing.  Firstly, let us assume that the local chord length and twist angle are defined at $N$ points along the wing; that is, at spanwise locations $y_i$, we know $c_i = c(y_i)$ and $\theta_i = \theta(y_i)$.
		
		$N - 1$ trailing-vortex lines attach to the wing at locations between these nodes, denoted by the index $j$.  The strength of each trailing vortex will be calculated as 
		
		\begin{equation}
			d\Gamma_j = \big (\Gamma_{i} - \Gamma_{i+1} \big ) \bigg |_{i = j} = \Gamma_{j} - \Gamma_{j+1}. \label{eq:dGamma}	
		\end{equation}
		That is, the strength of the $j$-th trailing vortex is the difference between the bound circulations at the adjacent nodes $i = j$ and $j+1$.
				
		The total downwash at the $i$-th node due to all the trailing vortex lines (numbered $j = 1$ to $N-1$) is
		
		\begin{equation}
			v_i = \sum_{j=1}^{N-1} \frac{d\Gamma_j}{4\pi(y_j - y_i)}
		\end{equation} 
		where $v_i$ is positive downwards and $d\Gamma_j$ is counter-clockwise-positive when viewed from downstream of the wing.  We can define an influence matrix as
		
		\begin{equation}
			Y_{ij} = \frac{1}{4\pi(y_j - y_i)}, \label{eq:Yij}
		\end{equation}
		such that 
		
		\begin{equation}
			v_i = \sum_{j=1}^N  Y_{ij} \,  d\Gamma_j = Y_{ij} \,  d\Gamma_j, \label{eq:vi}
		\end{equation}
		where repeated indices indicate summation.  In other words, $v_i$ is a vector that results from the multiplication of the matrix $Y_{ij}$ with the vector $d\Gamma_j$.
		
		To remain consistent with Helmholtz's vortex laws, we need to enforce $\Gamma_0$ and $\Gamma_N = 0$.  Accordingly, we need only solve for $\Gamma_i$ and $v_i$ for $i \in [2,N-1]$.
		
		Assuming that downwash is small relative to the freestream velocity (i.e., $v_i \ll U$), we can invoke the small-angle approximation to estimate the local effective angle of attack at each node as
		
		\begin{equation}
			\alpha_i = \theta_i + \Theta - \frac{v_i}{U}, \label{eq:alpha}
		\end{equation}
		where all angles are in radians.  
		
		To continue, we need to have a known or presumed relationship between $\Gamma$ and $\alpha$.  For a thin airfoil, the relationship is
		
		\begin{equation}
			\Gamma_i = \pi c U (\alpha_i - \alpha_0), \label{eq:Gamma}
		\end{equation}
		where $\alpha_0$ is the zero-lift angle of attack.  Recall that $\alpha_0 = 0$ for symmetric thin airfoils, but $\alpha_0 < 0$ when the wing is positively cambered (concave down, convex up).  If your airfoil is not thin, or if you want to account for viscous effects on lift, you can use an alternative relationship $\Gamma = \Gamma(\alpha)$: e.g., the results of an inviscid vortex panel method, Xfoil data, or the results of CFD simulations.  In any case, $\Gamma_i$ can be calculated or interpolated once $\alpha_i$ is known. 
		
		We now have enough information to solve, but we need to guess an initial circulation distribution to estimate the downwash distribution; then, we must calculate the circulation resulting from that downwash and iterate until the circulation converges.
		
		With a known circulation distribution, we calculate forces.  Here, we need to be deliberate with our language.  When we refer to local lift and local drag ($l$ and $d$), we are referring to the force per unit length acting normal and parallel to the local incident velocity (which depends on the freestream velocity and the downwash).  When we refer to total lift or drag on the wing ($L$ and $D$), we mean the force normal and parallel to the undisturbed freestream flow. Accordingly, we calculate 
		
		\begin{align}
			L &= \int_{-b/2}^{b/2} \bigg\{l(y) \cos\bigg(\frac{v(y)}{U}\bigg) - d(y) \sin\bigg(\frac{v(y)}{U}\bigg) \bigg\} dy, \\
			D &= \int_{-b/2}^{b/2} \bigg\{l(y) \sin\bigg(\frac{v(y)}{U}\bigg) + d(y) \cos\bigg(\frac{v(y)}{U}\bigg) \bigg\} dy, 
		\end{align}
		where $y=0$ is taken to lie at the midspan.  	
				
		To simplify the implementation in code, let us define the following dimensionless variables:
		
		\begin{align}
			&\Gamma_i^* = \frac{\Gamma_i}{b U}; 
			&c_i^* = \frac{c_i}{b}; &&y_i^* = \frac{y_i}{b}; &&v_i^* = \frac{v_i}{U}.
		\end{align}
			
		To these, we add the local and total force coefficients:
		
		\begin{align}
		&C_l = \frac{l}{\frac{1}{2} \rho c U^2}; &C_d = \frac{d}{\frac{1}{2} \rho c U^2}; &&C_L = \frac{L}{\frac{1}{2}\rho S U^2}; &&C_D = \frac{D}{\frac{1}{2}\rho S U^2},
		\end{align}
		where the total force coefficients are calculated from the local coefficients according to
		
		\begin{align}
			C_L &= \frac{b^2}{S} \int_{-1/2}^{1/2} c^* \bigg( c_l \cos v^* - c_d \sin v^* \bigg) dy^* \label{eq:CL} \\
			C_D &= \frac{b^2}{S} \int_{-1/2}^{1/2} c^* \bigg( c_l \sin v^* + c_d \cos v^* \bigg) dy^* \label{eq:CD}
		\end{align}
		
		At this point, you need to choose the model or the data you will use to determine local force coefficients.  Regardless of your choice, you will assume that the Kutta-Joukowsky theorem holds locally to relate $l = \rho U \Gamma$, from which we obtain 
		
		\begin{equation}
			\Gamma_i^* = \frac{1}{2} c_l c^*. \label{eq:Gammacl}
		\end{equation}
		
		Let us now finalize an algorithm that uses use thin-airfoil theory as the source of our local lift coefficients, and let us set $c_d = 0$, which amounts to an assumption of high Reynolds number and attached flow.  Using our dimensionless variables, our LLT algorithm can then be summarized as follows.
		
		\begin{callouttitle}{LLT Algorithm with Thin-Airfoil Theory}
			
			In this iterative algorithm, the index for iteration number will be $k$, appended to the iterated variable as a pre-superscript.
			
				\begin{enumerate}
					\item Guess values for $\prescript{k}{}{\Gamma_i^*}$, ensuring that $\Gamma_0^* = 0$ and $\Gamma_N^* = 0$.
					\item Calculate $d\Gamma_j^* = \Gamma_{i}^* - \Gamma_{i+1}^*$.
					\item Calculate the influence matrix, $Y_{ij}^* = \dfrac{1}{4\pi\,(y_j^* - y_i^*)}$.
					\item Calculate $v_i^* = Y^*_{ij} \, d\Gamma^*_j$ for $i \in [2,N-1]$.
					\item Calculate $\alpha_i = \theta_i + \Theta - v_i^*$ for $i \in [2,N-1]$.
					\item Calculate $\prescript{new}{}\Gamma_i^* = \pi c^*_i (\alpha_i - \alpha_0)$ for $i \in [2,N-1]$. \label{Gammanew}
					\item Calculate error and update your guess, $\prescript{k+1}{}\Gamma_i^*$. See note below for tips on doing so.
					\item Repeat the foregoing steps until $\Gamma_i^*$ converges.
					\item Calculate local lift coefficients as $c_{l,i} = \dfrac{2 \Gamma_i^*}{c_i^*}$.
					\item Calculate your total lift and drag coefficients using numerical approximations of the integrals in equations (\ref{eq:CL}) and (\ref{eq:CD}).
				\end{enumerate}
				
%			For numerical stability, consider choosing an initial guess that satisfies $\prescript{0}{}{\Gamma_0} = 0$ and $\prescript{0}{}{\Gamma_N}$.  
			
			Numerical stability can be improved, if necessary, by implementing a relaxation factor, $r$, when you update your guess.  That is,
			
			\begin{equation}
				\prescript{k+1}{}\Gamma_i^* = \prescript{k}{}\Gamma_i^* + r \, (\prescript{new}{}\Gamma_i^* - \prescript{k}{}\Gamma_i^*), \nonumber
			\end{equation}
			where $0< r < 1$ will improve stability over the simple choice of $\prescript{k+1}{}\Gamma_i^* = \prescript{new}{}\Gamma_i^*$, at the expense of requiring more iterations to converge.
			
			If you would like to use empirical airfoil data instead of thin-airfoil theory, replace step \ref{Gammanew} with an interpolation of empirical data that gives $c_l = c_l(\alpha)$ and use equation (\ref{eq:Gammacl}) to obtain the required relationship $\Gamma = \Gamma(\alpha)$.
				
				
				
		\end{callouttitle}
		
		
		
		

		

		
		
		
		
		


		
		% Bibliography at the very end
	\bibliographystyle{myplainnat}
%	\bibliography{allrefs}

	}
	
\end{document}
